\songtitle{Antiphona: De Atheismi Utilitate}

\tag*{2026}\tag{quora-post}

\begin{notes}
Lyrics by Grok based on a Quora post by Carlos Thompson
\end{notes}
\begin{style}
gregorian chant, liturgical antiphony, Latin text setting, solo cantor, responsorial choir, unison chant, solemn modal melody, slow 72 BPM, resonant stone chapel, natural reverb, sparse organ drone, plainchant cadence, gradual buildup, sectional contrast, austere medievalism, sacred dialogue, ritual gravity\\
new age track as background
\end{style}
\begin{lyrics}
\songpart{Introitus}
\annotation{choir} Utrum atheismus utilis est practicus?\\
\annotation{solo} Tam utilis quam capilli rufi naturales,\\
Vel statura centum septuaginta.\\
Quare utilis esse debet?\\
Condicio neutra, iudicium utilitatis vanum.

\songpart{Versus I}
Non confidere amico imaginario,\\
Utilis in circumstantiis quibusdam:\\
Problemata fronte directa solvis melius,\\
Non petens phantasma ut agat pro te.\\
Sed fides etiam iuvat:\\
In rebus extra potestatem,\\
Minus anxius si credis providentiam.

\songpart{Responsorium}
Credens tamen responsabilitatem sumere potest\\
Ut atheus;\\
Atheus vero stoica fortitudine\\
Adversitates incertas sustinet.

\songpart{Versus II (ex oratione serenitatis)}
Deus, dona mihi serenitatem\\
Accipere quae mutare non possum,\\
Fortitudinem mutare quae possum,\\
Sapientiam discernere differentiam.
— Reinhold Niebuhr.

\songpart{Responsorium Athei}
Non rogo Deum ut det mihi haec,\\
Sed ipse quaero serenitatem, fortitudinem, sapientiam.\\
Credens autem, spiritu orationis,\\
“Mutare quae possum” includit\\
Persequi haec dona, non passive exspectare.

\songpart{Conclusio (Antiphona finalis)}
Atheismus non est instrumentum,\\
Sed eventus fortuitus, accidens.\\
Quaerere utilitatem eius practicam\\
Error categoricus est.\\
Neutralis ut natura ipsa.\\
Amen.
\end{lyrics}

